\documentclass[12pt]{article}
\usepackage[utf8]{inputenc}
\usepackage[spanish]{babel}
\usepackage{amsmath, amssymb}
\usepackage{booktabs}
\usepackage{geometry}
\usepackage{tikz}
\usepackage{pgfplots}
\pgfplotsset{compat=1.18}
\usetikzlibrary{fillbetween}
\usepackage{float}
\geometry{margin=2cm}
\usepackage{hyperref}

\usepackage{subcaption}

\begin{document}

\vspace{5cm}



\begin{center}
\begin{large}
\textbf{\\Guía certamen 2: Econometría ICS-294}
\end{large}\\
\end{center}

\begin{enumerate}
   
\item 
En la tabla 1  se muestran los resultados de una investigación que busca identificar los determinantes de los salarios de los profesores (avgsal).  Utilizan las siguientes variables explicativas 

\begin{itemize}
   \item \texttt{lunch}: porcentaje de estudiantes elegibles para almuerzo gratuito.
  \item \texttt{enrol}: matrícula (número de estudiantes).
  \item \texttt{staff}: personal por cada 1000 estudiantes.
 
  \item \texttt{math4}: porcentaje que aprueba la prueba de matemáticas de 4.\textsuperscript{o} grado.
  \item \texttt{story4}: porcentaje que aprueba la prueba de lectura de 4.\textsuperscript{o} grado.
  \item \texttt{bs}: razón beneficios/salario, definida como \texttt{avgben/avgsal}.
 \texttt{lenrol}: logaritmo natural de \texttt{enrol}.
  \item \texttt{lstaff}: logaritmo natural de \texttt{staff}.
\end{itemize}

% Table created by stargazer v.5.2.3 by Marek Hlavac, Social Policy Institute. E-mail: marek.hlavac at gmail.com
% Date and time: Tue, Oct 14, 2025 - 11:46:19 AM
\begin{table}[H] \centering 
  \caption{Salarios de los profesores}  \label{tabla_reg}%\caption{Comparación: Regresiones} 

\begin{tabular}{@{\extracolsep{5pt}}lcc} 
\\[-1.8ex]\hline 
\hline \\[-1.8ex] 
 & \multicolumn{2}{c}{\textit{Dependent variable:}} \\ 
\cline{2-3} 
\\[-1.8ex] & \multicolumn{2}{c}{log(avgsal)} \\ 
\\[-1.8ex] & (1) & (2)\\ 
\hline \\[-1.8ex] 
 bs & $-$0.516$^{***}$ & $-$0.518$^{***}$ \\ 
  & (0.110) & (0.110) \\ 
  & & \\ 
 lenrol & $-$0.028$^{***}$ & $-$0.028$^{***}$ \\ 
  & (0.008) & (0.008) \\ 
  & & \\ 
 lstaff & $-$0.691$^{***}$ & $-$0.688$^{***}$ \\ 
  & (0.018) & (0.018) \\ 
  & & \\ 
 lunch & $-$0.001$^{***}$ & $-$0.001$^{***}$ \\ 
  & (0.0002) & (0.0002) \\ 
  & & \\ 
 math4 &  & 0.0004 \\ 
  &  & (0.0004) \\ 
  & & \\ 
 story4 &  & $$0.001 \\ 
  &  & (0.0005) \\ 
  & & \\ 
 Constant & 13.831$^{***}$ & 13.850$^{***}$ \\ 
  & (0.110) & (0.113) \\ 
  & & \\ 
\hline \\[-1.8ex] 
Observations & 1,848 & 1,848 \\ 
R$^{2}$ & 0.488 & 0.493 \\ 

\hline 
\hline \\[-1.8ex] 
\textit{Note:}  & \multicolumn{2}{r}{$^{*}$p$<$0.1; $^{**}$p$<$0.05; $^{***}$p$<$0.01} \\ 
\end{tabular} 
\end{table} 

Debiese incluir las variables \texttt{math4} y \texttt{story4} en la investigación? Sustente su respuesta mediante la formulación y explicación de un test de hipótesis que permita evaluar su relevancia estadística en el modelo ($
F_{0.05, (2, 1841)} \approx 3.00$).


{\color{blue} Se desea contrastar la hipótesis nula de que las dos variables eliminadas no aportan información estadísticamente significativa al modelo. Para ello, se emplea el test $F$ de restricciones lineales:

\[
H_0: \beta_{math4} =\beta_{story4} =0 
\]

\[
H_1: \text{Al menos uno es } \neq 0
\]

Los resultados estimados son los siguientes:

\[
R^2_{UR} = 0.493, \quad R^2_{R} = 0.488, \quad n = 1848, \quad k_{UR} = 6, \quad q = 2
\]

El estadístico de prueba se calcula como:

\[
F = \frac{(R^2_{UR} - R^2_{R}) / q}{(1 - R^2_{UR}) / (n - k_{UR} - 1)}
\]

Sustituyendo los valores:

\[
F = \frac{(0.493 - 0.488)/2}{(1 - 0.493)/(1848 - 6 - 1)} 
   = \frac{0.005 / 2}{0.507 / 1841}
   = \frac{0.0025}{0.0002754} \approx 9
\]

Por lo tanto, el estadístico $F$ sigue una distribución:

\[
F_{(q,\, n - k_{UR} - 1)} = F_{(2,\,1841)}
\]

El valor crítico al $5\%$ de significancia es aproximadamente:

\[
F_{0.05, (2, 1841)} \approx 3.00
\]

Dado que:

\[
F_{calc} = 9 > 3.00
\]

se \textbf{rechaza la hipótesis nula} $H_0$. 

\bigskip

\noindent
\textbf Existe evidencia estadísticamente significativa al 5\% de que las dos variables contribuyen de manera conjunta a explicar la variabilidad de la variable dependiente. Debiese incluir ambas variables en la investigación.
}


\item Se plantea el modelo:

$$votosA = \beta_0 + \beta_1 \cdot ln(\text{gasto A}) + \beta_2 \cdot ln(\text{gasto B}) + \beta_3 \cdot forpart + u$$

\begin{table}[H]
\centering
\caption{Resultados de la estimación del modelo de regresión}
\label{tab:regresion}
\begin{tabular}{lcccc}
\toprule
\textbf{Variable} & \textbf{Coeficiente} & \textbf{Error estándar} & \textbf{t-valor} & \textbf{Pr(>|t|)} \\
\midrule
Intercepto & 45.07893 & 3.92631 & 11.48 & $<$2e-16$^{***}$ \\
$\ln(\text{gasto A})$ & 6.08332 & 0.38215 & 15.92 & $<$2e-16$^{***}$ \\
$\ln(\text{gasto B})$ & -6.61542 & 0.37882 & -17.46 & $<$2e-16$^{***}$ \\
\textit{forpart} & 0.15196 & 0.06202 & 2.45 & 0.0153$^{*}$ \\
\midrule
\multicolumn{5}{l}{\textit{Signif. codes:} $^{***}p<0.001$, $^{**}p<0.01$, $^{*}p<0.05$} \\
\midrule
\multicolumn{5}{l}{Error estándar residual: 7.712 (169 g.l.)} \\
\multicolumn{5}{l}{$R^2$ = 0.7926 \quad $R^2_{ajustado}$ = 0.7889} \\
\multicolumn{5}{l}{Estadístico F = 215.2 (3; 169) \quad p-valor $<$ 2.2e-16} \\
\bottomrule
\end{tabular}
\end{table}

Donde $ln(\text{gasto A})$ y $ln(\text{gasto B})$ son el logaritmo natural del gasto en campaña del candidato A y B, respectivamente, $forpart$ es una medida de la fortaleza del partido del candidato A (porcentaje de votos que ovtuvo el partido A en la elección presidencial más reciente), y votosA el porcentaje de votos recibidos por el candidato A.\\

Además, se tiene que la matriz de varianzas-covarianzas de los estimadores es:

\begin{table}[H]
    \centering
    \begin{tabular}{c|c|c|c|c}
 & (Intercept) & ln(gasto A) & ln(gasto B)	& forpart\\\hline
(Intercept)	& 15.4158713 & -0.394937598	& -0.874063434 & -0.176167017\\
ln(gasto A)	& -0.3949376 & 0.146038608 & -0.002681489 & -0.006546373\\
ln(gasto B)	& -0.8740634 & -0.002681489	& 0.143504821 & 0.003577341\\
forpart	& -0.1761670 & -0.006546373	& 0.003577341 & 0.003846245\\
    \end{tabular}
    %\caption{Caption}
    %\label{tab:placeholder}
\end{table}

Plantee el test de hipótesis de que un aumento de un 1\% en los gastos de A es compensado por un aumento de 1\% en los gastos de B, determine la veracidad de lo anterior al 5\% y explique.\\
\color{blue}
\textbf{Respuesta:} El test de hipótesis planteado es el siguiente: 

\begin{align*}
    H_0 &: \beta_1 = -\beta_2\\
    H_1 &: \beta_1 \neq -\beta_2
\end{align*}

\begin{align*}
    H_0 &: \beta_1 +\beta_2=0\\
H_1 &: \beta_1 +\beta_2 \neq 0 
\end{align*}
Luego, el cálculo de $t$ es:
$$t = \frac{\widehat{\beta}_1 + \widehat{\beta}_2}{\sqrt{\hat{V}(\widehat{\beta}_1 + \widehat{\beta}_2)}}$$

Reemplazando y aplicando la propiedad de la varianza de una suma\footnote{Los valores finales están calculados con todos los decimales, puede haber diferencias por aproximación}:

$$t = \frac{6.083 - 6.615}{\sqrt{0.146 + 0.144 + 2\cdot (-0.003)}} \approx \frac{0.53324}{−0.53210 }= −0.9978$$

Luego, el $|t_{calc}| < 1.96$, por lo tanto, no existe evidencia suficiente para rechazar la hipótesis nula y, efectivamente, un aumento de 1\% en el gasto en campaña del candidato A se ve contrarrestado por un aumento de 1\% en gasto de campaña del candidato B.
    \end{enumerate}


\end{enumerate}
    \end{document}